  
%
% $Id: slides.tex 4228 2006-06-21 21:55:12Z jjamor $
%
%
% Compilar a .pdf con LaTeX (pdflatex)
% Es necesario instalar Beamer (paquete latex-beamer en Debian)
%



\documentclass{beamer}
\usetheme{Warsaw}
\usebackgroundtemplate{\includegraphics[width=\paperwidth]{format/libresoft-bg.png}}
\usepackage[english]{babel}
\usepackage[latin1]{inputenc}
\usepackage{graphics}
\usepackage{amssymb} % Simbolos matematicos
\usepackage{url}
\usepackage{listings}
%\usepackage{html}
%\usepackage{hthtml}

%\definecolor{libresoftgreen}{RGB}{162,190,43}
%\definecolor{libresoftblue}{RGB}{0,98,143}

%\setbeamercolor{titlelike}{bg=libresoftgreen}

% for xml code. Usage: \begin{xml} xml code... \end{xml}
\lstnewenvironment{xml}
{\lstset{language=xml,numbers=left,stepnumber=2}}
{}


%% Metadatos del PDF.
\hypersetup{
  pdftitle={Layouts},
  pdfauthor={Roberto Calvo Palomino},
  pdfcreator={GSyC/Libresoft},
  pdfproducer=PDFLaTeX,
  pdfsubject={nn},
}
%%


\AtBeginSection[]
{
  \begin{frame}<presentation>
    \frametitle{Index}
    \tableofcontents[current]
  \end{frame}
}


\begin{document}

\title{Curso de Desarrollo en Android }
\subtitle{Layouts}
\institute{\\rocapal@libresoft.es\\
GSyC/Libresoft}
\author{Roberto Calvo Palomino}
\date{\today}

\frame{
\maketitle
\begin{center}
\includegraphics[width=6cm]{format/gsyc-urjc}
\hspace{0.5cm}
\includegraphics[width=2.5cm]{format/eoi}
\end{center}
}


%\title{Titulo corto}
%\author{Autores abreviado}


%% LICENCIA DE REDISTRIBUCION DE LAS TRANSPAS
\frame{
~
\vspace{4cm}

\begin{flushright}
{\tiny
(cc) 2011 Roberto Calvo Palomino. \\
Some rights reserved. This document is distributed under the Creative \\
            Commons Attribution-ShareAlike 3.0 licence, available in \\
            http://creativecommons.org/licenses/by-sa/3.0/


}
\end{flushright}
}

\section{Nomenglatura}
\begin{frame}
\frametitle{Medici�n en Pantalla}
\begin{itemize}
\item px (p�xeles)
\item dip (p�xeles independientes del dispositivo)
\item sp (p�xeles escalados, mejor para tama�os de texto)
\item pt (puntos)
\item in (inches)
\item mm (mil�metros)
\end{itemize}
\end{frame}

\begin{frame}
\frametitle{Diferentes Pantalla}
\begin{itemize}
\item Tama�o
\item Densidad (dependiente o no del pixel)
\item Aspecto de Ratio
\item Versi�n de plataforma
\end{itemize}

\begin{center}
\htmladdnormallink{Screens Support (Link)}{http://developer.android.com/guide/practices/screens\_support.html\#range}
\end{center}
\end{frame}

\section{Tipos de Layouts}

\begin{frame}
\frametitle{Introduccion}
\begin{itemize}
\item Unidad elemental de la interfaz gr�fica.
\item Contenedores de widgets que ocupan la pantalla
\item Relaci�n fuerte entre layouts para conseguir la interfaz deseada
\item fill\_parent: Ocupa la anchura o altura m�ximo que le deje su padre
\item wrap\_content: Ocupa la anchura o altura m�nima necesaria
\end{itemize}
\end{frame}

\begin{frame}
\frametitle{FrameLayout}
\begin{itemize}
\item Est� dise�ado para ocupar un area de la pantalla para mostrar un
  solo elemento
\item Puedes a�adir muchos hijos a este FrameLayout, pero todos los
  hijos aparecer�n alineados arriba a la izquierda de este contenedor. 
\item En este grupo de vistas no se puede especificar la posici�n de
  los hijos
\item No son los layouts m�s utilizados
\end{itemize}
\end{frame}


\begin{tiny}
\begin{block}{}
\begin{xml}
<?xml version=''1.0'' encoding=''utf-8''?>
<FrameLayout
xmlns:android=''http://schemas.android.com/apk/res/android''
    android:layout_width=''150dip''
    android:layout_height=''150dip''
    android:background=''#FF00FF00''
    >     
  <Button
    android:layout_width=''wrap_content'' 
    android:layout_height=''wrap_content'' 
    android:text=''Bot�n''
    />    
  <TextView
    android:layout_width=''wrap_content''
    android:layout_height=''wrap_content''
    android:text=''Hola que tal como est�s. Este texto aparece encima
    del bot�n''
  />    
</FrameLayout>
\end{xml}
\end{block}
\end{tiny}

\begin{frame}
\frametitle{LinearLayout}
\begin{itemize}
\item Dispone sus hijos en una sola columna o en una sola fila (configurable)
\item La direcci�n de cada fila puede ser especificada mediante
  \emph{setOrientation()} en c�digo o mediante
  \emph{android:orientation} en el XML. Por defecto, es horizontal
\item Tambi�n puedes especificar gravedad, la cual indica la
  alineaci�n de todos los elementos hijos
\item Soporta asignar peso a los hijos de forma
  individual. Este valor permite a los hijos que se expandan hasta
  llenar cualquier espacio libre de la pantalla.
\end{itemize}
\end{frame}

\begin{tiny}
\begin{block}{}
\begin{xml}
<?xml version=''1.0'' encoding=''utf-8''?>
<LinearLayout
xmlns:android=''http://schemas.android.com/apk/res/android''
    android:orientation=''horizontal''
    android:layout_width=''fill_parent''
    android:layout_height=''wrap_content''
    android:background=''#FF00FF00''
    >     
<TextView
    android:layout_width=''wrap_content''
    android:layout_height=''wrap_content''
    android:text=''Texto prueba''
    android:background=''#FFFFFFFF''
    android:layout_weight=''5''
    />
<Button
    android:layout_width=''wrap_content''
    android:layout_height=''wrap_content''
    android:text=''Bot�n''
    android:layout_weight=''1''
    />
</LinearLayout>
\end{xml}
\end{block}
\end{tiny}

\begin{frame}
\frametitle{RelativeLayout}
\begin{itemize}
\item El m�s flexible de todos los tipos de layouts. Puedes construir
  casi cualquier cosa con ellos.
\item La posici�n de los hijos son relativas al padre e incluso a
  otros hijos.
\item Nos permite realizar interfaces complejas y portables a
  diferentes tama�os de pantalla.
\end{itemize}
\end{frame}

\begin{tiny}
\begin{xml}
<?xml version=''1.0'' encoding=''utf-8''?>
<RelativeLayout
xmlns:android=''http://schemas.android.com/apk/res/android''
    android:layout_width=''fill_parent''
    android:layout_height=''fill_parent''
    >
    <TextView  
        android:id=''@+id/hello''
        android:layout_width=''fill_parent'' 
        android:layout_height=''wrap_content'' 
        android:text=''@string/hello''
        />
    <TextView
        android:id=''@+id/black_on_red''
        android:layout_width=''fill_parent''
        android:layout_height=''wrap_content''
        android:layout_below=''@id/hello''
        android:background=''#ff0000''
        android:textColor=''#000000''
        android:text=''black on red''
        />
    <TextView
        android:id=''@+id/update_me''
        android:layout_width=''fill_parent''
        android:layout_height=''wrap_content''
        android:layout_below=''@id/black_on_red''
        android:background=''#00ff00''
        android:textColor=''#000000''
        android:text=''update Text''
        /> 
    <TextView
        android:id=''@+id/left_wrap''
        android:layout_width=''wrap_content''
        android:layout_height=''wrap_content''
        android:layout_below=''@id/update_me''
        android:background=''#0000ff''
        android:textColor=''#000000''
        android:text=''wrap''
        />
    <TextView
        android:id=''@+id/right_fill''
        android:layout_width=''fill_parent''
        android:layout_height=''wrap_content''
        android:layout_below=''@id/update_me''
        android:layout_toRightOf=''@id/left_wrap''
        android:background=''#ff0000''
        android:textColor=''#000000''
        android:text=''fill''
    />
</RelativeLayout>
\end{xml}
\end{tiny}

\begin{frame}
\frametitle{TableLayout}
\begin{itemize}
\item Layout espec�fico para mostrar tablas
\item Se puede realizar tambi�n con RelativeLayouts
\item Columas y Filas son sus principales propiedades
\end{itemize}
\end{frame}


\begin{tiny}
\begin{xml}
<?xml version=''1.0'' encoding=''utf-8''?>
<TableLayout xmlns:android=''http://schemas.android.com/apk/res/android''
    android:layout_width=''wrap_content''
    android:layout_height=''fill_parent''
    android:stretchColumns=''1''>
    <TableRow>
        <TextView
            android:text=''prueba''
            android:padding=''3dip'' />
        <TextView
            android:text=''prueba2''            
            android:padding=''3dip'' />
    </TableRow>

    <TableRow>
        <TextView
            android:text=''prueba3''
            android:padding=''3dip'' />
        <TextView
            android:text=''prueba4''
            android:padding=''3dip'' />
    </TableRow>
</TableLayout>
\end{xml}
\end{tiny}



\section{Depuraci�n de Vistas}
\begin{frame}
\frametitle{HierarchyViewer y LayoutOpt}
\begin{itemize}
\item Herramientas para depurar nuestras interfaces
\item �til para poder optimizar recursos en nuestras aplicaciones
\item Se encuentra en \bf{\$DIR\_SDK\_ANDROID/tools/}
\end{itemize}

\begin{center}
  \htmladdnormallink{Depurar Interfaces (Link)}{http://developer.android.com/guide/developing/debugging/debugging-ui.html\#layoutopt}
\end{center}

\end{frame}

\begin{frame}
\frametitle{Optimizaci�n}
\begin{itemize}
\item Nunca usar AbsoluteLayout
\item Usar el m�nimo n�mero de vistas posibles
\item Nunca m�s de 80 layouts en una misma vista, y nunca m�s de 10
  layouts anidados
\item Pantallas y vistas sencillas
\item Utilizar recursos apropiados para cada tipo de pantalla
\end{itemize}
\end{frame}
\end{document}
